% FILE: artifacts/_common/style/figure-toolkit.tex
% ==================================================================================================
% THE FIGURE TOOLKIT of the house style: the shared 3-D camera, the TikZ styles of the drawn objects
% (arrows, vectors, drop lines, span planes, the neutral in-figure shadow), the coordinate dots, the
% tapered arrow and the connector. A document inputs it after the colours
% (artifacts/_common/style/colors.tex), whose tokens the styles name.

% ===== THE SHARED 3-D CAMERA ======================================================================
% ONE elevation / azimuth for every tikz-3dplot figure that is NOT the mechanism tree or the supplement
% (those keep their own \mtcamtheta / \spcamtheta because they need a per-mode azimuth). \geometryplot,
% \inputplane / \responsedecomp and the two activity-space plots all read these, so the whole poster
% projects one scene from one viewpoint and a camera retune is a ONE-LINE edit here.
%   theta = ELEVATION (polar angle from the vertical z = n1). At 45 the view is OBLIQUE (the archive's
%   attitude, matching the mechanism tree): a tilted plane reads as tilted, not as a floor seen from above
%   (which is what the earlier low 29 gave). The span planes are kept OPEN not by lowering the camera but
%   by GENTLING the generators (z-share ~0.24): a plane whose normal sits only ~19 deg off vertical is
%   ~45% open even at this oblique view, while its residual still rises. phi = AZIMUTH, fixed so world-z
%   projects up, y right, x down-left (n1 up / n2 right / n3 receding), the archive's orientation.
\def\camtheta{57}\def\camphi{118}
% The ANSWER FACETS (\geometryplot) view the scene from FURTHER ABOVE (a smaller theta): this droops the
% floor axes more to the south, so the components plane reads as clearly distinct from the horizontal
% floor. Its plane is steepened to compensate (geometries.tex) so the plane's projected surface still
% matches the approach plots' look. phi is shared.
\def\facetcamtheta{38}
% ===== SHARED TikZ STYLES =========================================================================
% THE ONLY COPY OF EACH OF THESE NUMBERS. A figure that needs one of these devices takes the style and
% overrides only what is genuinely different (colour, and -- where noted -- weight); it never re-spells
% the arrow head or the shadow parameters. Retune here, or nowhere.
\tikzset{
  % ---- shadows -------------------------------------------------------------------------------------
  % Shallow, neutral depth for objects INSIDE a figure (tree strips, option pills, taxonomy cells, leaf
  % plates). Lighter than `posterglow': they must not read as cards floating over the card holding them.
  softshadow/.style={blur shadow={shadow xshift=0pt, shadow yshift=-1.8mm, shadow blur radius=3.6mm,
    shadow blur steps=9, shadow opacity=34}},
  % ---- arrows / connectors ------------------------------------------------------------------------
  flowarrow/.style={-{Stealth[length=19pt,width=15pt]}, line width=2.2pt,
    rounded corners=8pt},
  % THE FRAME AXIS -- one weight and one head for every coordinate frame on the poster (the neuron
  % triples of \inputplane / \responsedecomp / \geometryplot, and the mechanism tree's intrinsic triple).
  % At 1.1 pt / 8 pt a frame prints as a hairline behind 3.4 pt vectors with invisible heads, so the
  % axes read as construction lines. The COLOUR is not here: each figure supplies its own `borderfunc' tint.
  frameaxis/.style={-{Stealth[length=15pt, width=11pt]}, line width=2.2pt},
  % A NAMED VECTOR -- the objects a figure is about. \geometryplot draws them at the declared 4 pt;
  % the Approach's two probe plots pass `line width=3.6pt' because they carry more arrows in less room.
  namedvec/.style={-{Stealth[length=16pt, width=12pt]}, line width=4pt},
  % THE ORTHOGONAL RESIDUAL (eta_perp): solid brown, one weight across the geometry figures.
  orthvec/.style={-{Stealth[length=13pt, width=10pt]}, line width=2.8pt, draw=outplane},
  % THE BIAS ARC -- the in-plane rotation off the linear sum, drawn concentric with it.
  biasarc/.style={-{Stealth[length=8pt]}, line width=1.4pt, draw=outplaneannot},
  % (The coordinate mark / dot lives further down as \dotonvec / \dotonaxis / \dotonplane, sized as a
  % fraction of each figure's axis length; see the block below the tikzset.)
  % THE LINEAR SUM (the proportional reference), one convention document-wide. Every geometry is
  % measured against it -- bias IS the deviation from it -- so it must be recognisable as the SAME
  % object in the taxonomy cells and in the tree. Absolute values, so `scale' cannot change its look.
  refvec/.style={-{Stealth[length=13pt, width=10pt]}, line width=2.8pt, draw=textsecondary,
    dash pattern=on 6pt off 3pt},
  % THE DROP LINE from a vector's tip (or its shadow) to the axis a coordinate is read on. ONE dash
  % geometry across both documents, so a reader crossing from poster to supplement meets the same mark.
  % Only the WEIGHT differs by figure family (`coorddrop' large, `coorddropsm' page-scale), because line
  % widths do not scale with a tikz `scale'. Override the colour with draw=.
  coorddrop/.style={line width=1.5pt, draw=texttertiary, opacity=0.9,
    dash pattern=on 3pt off 3pt, line cap=round},
  coorddropsm/.style={line width=0.9pt, draw=texttertiary, opacity=0.9,
    dash pattern=on 2.4pt off 2.4pt, line cap=round},
  % ---- the components plane -----------------------------------------------------------------------
  % THE HOUSE RULE FOR EVERY SPAN PLANE ON THE POSTER, stated once here: the plane IS the parallelogram
  % of its two generators -- 0 -> a -> a+b -> b -- so each arrow ENDS ON THE CORNER of the edge it
  % defines, and an inflated copy leaves every arrowhead floating short of its own corner. Anything
  % drawn inside the patch must keep BOTH coefficients < 1; the one licensed exception is
  % \geometryplot's unbiased cell, where the mixture IS the exact sum k1+k2 and lands on the far corner.
  % `compplane' IS "the plane of the components" here; two planes are told apart by ATTITUDE, never by
  % hue. Weight is overridden where a figure is drawn at a smaller scale.
  spanfill/.style={compplane, opacity=0.32},
  spanedge/.style={draw=compplane!75, line width=0.8pt},
}
% THE SHARED COORDINATE DOT -- SIZED FROM THE LINE WIDTH of the line it sits on, in THREE kinds, one per
% kind of line the project draws:
%   \dotonvec   -- on a RESPONSE / component vector (the thick rays)
%   \dotonaxis  -- on a FRAME AXIS or basis vector (incl. eigen-axes and eigenvectors)
%   \dotonplane -- a foot on a PLANE, marked by the dotted drop lines
% A dot reads as a bead threaded on its line, so its size follows THAT LINE'S WIDTH: a thick response
% vector carries a big bead, a hair-thin panel axis a small one. Because a figure drawn small also draws
% THINNER lines, this tracks figure size for free, with no per-figure scale knob. The law is affine --
% diameter = \dotbase + \dotgain x width -- not a bare multiple: the small \dotbase keeps a bead on a
% 0.7pt supplement axis legible instead of shrinking it to a speck, while \dotgain sets how fast the bead
% grows with the line. (Calibrated so the poster's 3.6pt response vectors carry the archive's ~8pt bead.)
% Each figure \def's only the widths it uses -- \lwvec / \lwaxis / \lwplane -- to its OWN line weights.
%   DRAWN AS NODES, not `circle' paths: inside `tdplot_main_coords' a `circle' is projected to an ELLIPSE,
%   whereas a circle NODE is placed at the projected point and drawn round in the page. No `transform
%   shape' (the size is absolute pt, like a line width). White halo under a coloured disc; halo/disc =
%   \dothaloratio, the rim being the gap between the two filled discs.
\def\dotbase{1.5}\def\dotgain{1.83}\def\dothaloratio{1.32}   % diameter (pt) = base + gain*linewidth
\def\lwvec{3.6}\def\lwaxis{2.0}\def\lwplane{1.5}             % defaults; each figure sets its actual weights
\newcommand{\dotcore}[3]{% #1 line width (pt number)  #2 fill colour  #3 (coordinate)
  \node[circle, fill=white, minimum size=\fpeval{\dothaloratio*(\dotbase+\dotgain*(#1))}pt, inner sep=0pt] at #3 {};%
  \node[circle, fill=#2,    minimum size=\fpeval{\dotbase+\dotgain*(#1)}pt,                 inner sep=0pt] at #3 {};}
\newcommand{\dotonvec}[2]{\dotcore{\lwvec}{#1}{#2}}
\newcommand{\dotonaxis}[2]{\dotcore{\lwaxis}{#1}{#2}}
\newcommand{\dotonplane}[2]{\dotcore{\lwplane}{#1}{#2}}

% ===== TAPERED / CURVED CONSEQUENCE ARROW =========================================================
% A FILLED, tapered, curved arrow (a narrow tail swelling to a wide head with a triangular point) --
% the "swoosh" connector between panels. No standard package draws a tapered CURVED body, so it is a
% filled bezier ribbon: the spine is a quadratic bend, the two edges are the spine offset by a width
% that grows from tail to neck, and a triangle caps the head. Widths are ABSOLUTE (pt), so the arrow
% keeps its weight regardless of the picture `scale' (like a line width or an arrowhead).
%   \taperarrow[<fill>][<extra path opts>]{<tail>}{<head>}{<bend>}{<tail hw>}{<neck hw>}{<ah len>}{<ah hw>}
%   bend = signed perpendicular offset of the spine's midpoint (0 = straight; +/- curves the swoosh).
%   The SECOND optional argument is a raw TikZ option list applied to the ribbon path together with the
%   fill -- pass `draw=<colour>, line width=<w>' for a border and a shadow key (e.g. `blur shadow={...}')
%   for a drop shadow. Empty by default, so bare `\taperarrow[<colour>]{...}' calls are unchanged.
\NewDocumentCommand{\taperarrow}{O{black} O{} m m m m m m m}{%
\begin{scope}
  \path
    let \p1 = ($#4-#3$), \n1 = {veclen(\x1,\y1)},
        \p2 = (\x1/\n1, \y1/\n1),            % unit direction (tail -> head)
        \p3 = (-\y1/\n1, \x1/\n1)            % unit perpendicular
    in
    coordinate (T@)  at #3
    coordinate (H@)  at #4
    coordinate (N@)  at ($#4-#8*(\p2)$)                      % neck (base of the arrowhead)
    coordinate (Bc@) at ($($#3!0.5!(N@)$)+#5*(\p3)$)         % bent spine mid
    coordinate (Tt@) at ($#3+#6*(\p3)$)  coordinate (Tb@) at ($#3-#6*(\p3)$)
    coordinate (Nt@) at ($(N@)+#7*(\p3)$) coordinate (Nb@) at ($(N@)-#7*(\p3)$)
    % control widths = the MEAN of the tail and neck half-widths, not the neck's: at the mean the four
    % control points of each edge are collinear when bend=0, so the taper is a straight ramp instead of
    % a curve that balloons out near the tail (visible only on a long arrow; short ones barely bent).
    coordinate (Ct@) at ($(Bc@)+0.5*#6*(\p3)+0.5*#7*(\p3)$)
    coordinate (Cb@) at ($(Bc@)-0.5*#6*(\p3)-0.5*#7*(\p3)$)
    coordinate (Aw1@) at ($(N@)+#9*(\p3)$) coordinate (Aw2@) at ($(N@)-#9*(\p3)$);
  \path[fill=#1, #2]
    (Tt@) .. controls (Ct@) and (Ct@) .. (Nt@) --
    (Aw1@) -- (H@) -- (Aw2@) --
    (Nb@) .. controls (Cb@) and (Cb@) .. (Tb@) --
    cycle;
\end{scope}}

% ===== THE POSTER'S CONNECTOR =====================================================================
% \connarrow{<tail>}{<tip>} -- ONE arrow shape for ONE relation ("leads to"), and the only \taperarrow
% call site the poster has. Every connector on the page runs through it: the two premise arrows into
% the medallion and the four question -> approach -> answer arrows down the columns. A second styling
% would assert a second kind of relation, which is exactly what the label system forbids.
% THE GEOMETRY IS TUNED FOR A SHORT ARROW. These runs are only ~2.6-3.7 cm, so the neck is thick and
% the head wide; a slender arrow at that length reads as a scratch. Straight (bend 0) because every
% run is axis-aligned.
% The border and shadow options go in \taperarrow's SECOND optional argument, and they are passed
% LITERALLY rather than through a macro -- pgfkeys must see the commas to split the list.
\newcommand{\connarrow}[2]{%
  \taperarrow[spinefill!45][draw=spineedge, line width=1.8pt, softshadow]%
    {#1}{#2}{0pt}{2pt}{18pt}{22pt}{44pt}}
% How far each end of a connector reaches INTO the two objects it joins, so it reads as one continuous
% link rather than as a floating glyph in a gap.
\def\connov{0.5}
